\subsection{Complete calculation of a mulch benefit}
\exbox{Mulch, tillage and the feedback loop}{Compare the evaporation under bare soil, straw mulch (28 per cent reduction) and plastic film (55 per cent) over a 10-day spell, find the length of stage-1 drying, the change of $K_s$ from tillage, and the feedback on the irrigation interval.}\sol
\ans Evaporation from bare soil: potential $E_0=4.5$ mm\,d$^{-1}$, mulch reduction factor $0.00$, so $E=E_0(1-k)=4.50$ mm\,d$^{-1}$; over a 10-day drying spell $45.0$ mm, saving $0.0$ mm
\ans Evaporation from straw mulch 5 t/ha: potential $E_0=4.5$ mm\,d$^{-1}$, mulch reduction factor $0.28$, so $E=E_0(1-k)=3.24$ mm\,d$^{-1}$; over a 10-day drying spell $32.4$ mm, saving $12.6$ mm
\ans Evaporation from plastic film: potential $E_0=4.5$ mm\,d$^{-1}$, mulch reduction factor $0.55$, so $E=E_0(1-k)=2.02$ mm\,d$^{-1}$; over a 10-day drying spell $20.2$ mm, saving $24.8$ mm
\ans Stage-1 evaporation (Ritchie) ends when $U=REW=9$ mm; with $E=4.5$ it lasts $9/4.5=2.0$ d on bare soil and $9/3.24=2.8$ d under straw
\ans Tillage: raising the porosity from 0.43 to 0.49 raises $K_s$ by the Kozeny--Carman ratio \[\left[\frac{0.49^3/(1-0.49)^2}{0.43^3/(1-0.43)^2}\right]=1.85\]
\ans Feedback loop: lower evaporation $\to$ wetter surface $\to$ longer interval; with a mulch saving of $12.6$ mm per 10 days the interval of 7 days becomes 8.3 days for the same TAW \finans


